\chapter{Dental Hygiene}
\index{Dental Hygiene}

\section*{Definition and Overview}
Dental hygiene refers to the daily practices that keep the teeth, gums, and mouth clean and healthy, including brushing, cleaning between the teeth, and using fluoride. Good dental hygiene prevents tooth decay and gum disease, which are among the most common health problems worldwide and are largely preventable. Effective hygiene removes plaque, the sticky film of bacteria that forms constantly on the teeth, before it can cause damage. This chapter provides practical, evidence-based guidance on brushing technique, interdental cleaning, and other habits that protect oral health.

\section*{Epidemiology}
Most people brush their teeth at least once daily, but fewer clean between their teeth, and gum disease remains common. Around one in three adults has moderate to severe gum disease, and tooth decay affects most people at some point. Poor dental hygiene is a major driver of these conditions, and it is more common in disadvantaged groups, older adults, and people with disabilities. Regular professional dental care combined with good home hygiene is associated with dramatically lower rates of decay, gum disease, and tooth loss.

\section*{Aetiology and Pathophysiology}
\begin{itemize}
    \item \textbf{Plaque formation:} bacteria, food debris, and saliva form a sticky film on teeth within hours of cleaning
    \item \textbf{Acid production:} plaque bacteria ferment sugars, producing acid that dissolves enamel and causes decay
    \item \textbf{Gum inflammation:} plaque at the gum line inflames the gums (gingivitis), which can progress to periodontitis with bone loss
    \item \textbf{Calculus:} plaque that is not removed hardens into tartar, which cannot be removed by brushing and provides a rough surface for more plaque
    \item \textbf{Fluoride:} strengthens enamel, reverses early decay, and reduces acid production
    \item \textbf{Prevention mechanism:} regular removal of plaque keeps acid exposure low and gums healthy
\end{itemize}

\section*{Clinical Features}
\begin{itemize}
    \item Healthy gums are pink, firm, and do not bleed on brushing
    \item Signs of poor hygiene: bleeding gums, bad breath, a fuzzy feeling on the teeth, and visible plaque or tartar
    \item Early gum disease is painless, so daily hygiene and professional checks matter
    \item Tooth sensitivity and white or brown spots may indicate early decay
    \item Effective hygiene results in smooth, clean teeth and fresh breath
\end{itemize}

\section*{Differential Diagnosis}
\begin{itemize}
    \item Gingivitis from plaque versus gum disease from other causes, such as medication effects or hormonal changes
    \item Bad breath from plaque versus from tonsil stones, sinus disease, or reflux
    \item Early decay versus staining
    \item Tooth sensitivity from gum recession versus from erosion or decay
\end{itemize}

\section*{When to Seek Medical Care}
Everyone should have a dental check-up at least once a year, or as recommended. See a dentist or dental hygienist promptly for bleeding gums, persistent bad breath, sensitivity, or visible plaque and tartar that will not brush away.

\textbf{Contact your local emergency services for:}
\begin{itemize}
    \item Severe pain, swelling, or signs of dental infection
    \item Uncontrolled bleeding from the mouth
    \item Serious dental trauma
\end{itemize}

\section*{Diagnosis and Investigations}
\begin{itemize}
    \item \textbf{Dental examination:} assessment of plaque levels, gum health, and decay
    \item \textbf{Periodontal probing:} measurement of gum pockets and bleeding to assess gum disease
    \item \textbf{Disclosing tablets:} colour plaque so that missed areas become visible
    \item \textbf{Risk assessment:} review of hygiene habits, diet, and medical history
    \item \textbf{Professional cleaning:} removal of plaque and tartar, and polishing
\end{itemize}

\section*{Management}
\begin{itemize}
    \item \textbf{Brushing:} brush twice daily for two minutes with fluoride toothpaste, using gentle circular motions
    \item \textbf{Interdental cleaning:} floss or interdental brushes once daily to clean between the teeth where brushing cannot reach
    \item \textbf{Professional care:} regular cleaning by a dentist or hygienist removes tartar
    \item \textbf{Fluoride:} use fluoride toothpaste and consider mouthwash for high-risk individuals
    \item \textbf{Electric toothbrushes:} are often more effective at removing plaque than manual brushes
    \item \textbf{Tongue cleaning:} gently brush or scrape the tongue to reduce bad breath and bacteria
    \item \textbf{Treatment of gum disease:} scaling and root planing when periodontitis is present
\end{itemize}

\section*{Complications}
\begin{itemize}
    \item Gingivitis progressing to periodontitis, with gum recession, bone loss, and loose teeth
    \item Tooth decay leading to fillings, root canal treatment, or extraction
    \item Bad breath and staining affecting confidence
    \item Tooth loss and its effects on eating and appearance
    \item Links between gum disease and systemic conditions such as diabetes and heart disease
    \item Cost of treating advanced disease compared with prevention
\end{itemize}

\section*{Prognosis and Outcomes}
With consistent, effective dental hygiene, decay and gum disease are largely preventable, and most people can keep their natural teeth for life. Gingivitis is fully reversible with improved cleaning and professional treatment, and early periodontitis can be halted. The benefits of good dental hygiene extend to general health and quality of life, and the long-term outlook is excellent for those who maintain the habits.

\section*{Prevention and Self-Care}
\begin{itemize}
    \item Brush twice daily for two minutes with fluoride toothpaste
    \item Floss or use interdental brushes daily
    \item Replace the toothbrush every three months or after illness
    \item Limit sugary and acidic foods and drinks
    \item Drink fluoridated tap water
    \item Visit the dentist regularly
    \item Do not smoke, and limit alcohol
\end{itemize}

\section*{Sources and Further Reading}
The following reputable sources provide further information for healthcare students and patients seeking reliable, current advice about dental hygiene.

\begin{itemize}
    \item Australian Dental Association. \textit{Brushing and flossing: how to look after your teeth.} https://www.ada.org.au/
    \item NHS. \textit{Teeth cleaning: how to brush and floss properly.} https://www.nhs.uk/live-well/healthy-teeth-and-gums/
    \item Australian Dental and Oral Health Therapists Association. \textit{Dental hygiene services in Australia.} https://www.adahta.com.au/
    \item World Health Organization. \textit{Oral health: key facts.} https://www.who.int/
    \item Healthdirect Australia. \textit{How to clean your teeth properly.} https://www.healthdirect.gov.au/
    \item MedlinePlus. \textit{Teeth and gum care.} https://medlineplus.gov/
\end{itemize}
